\reviewsection{The Environment Determines Which Competence Becomes Visible}

This question is personal for me as a Latino Dominican, bilingual, neurodivergent, and chronically ill learner and educator. I know how quickly an access barrier can be rewritten as a character judgment. Fatigue can become lack of commitment. Processing time can become lack of knowledge. Language difference can become confusion. A quieter or delayed response can become lack of interest. Meaningful communication requires adults to examine the conditions under which competence is being judged. In my K--5 technology teaching placement, I saw participation change when I reduced language demands, modeled the exact technology steps, made the sequence visible, and allowed familiar tools to remain available. The learner's competence had not suddenly appeared; the conditions for showing it had changed. In one documented interaction, a Spanish-language interface did not stop progress because visual demonstration, pointing, and short directions carried the instruction while the student completed the task \citep{placementEvidence2026}. Language access and partner behavior were not extras around learning. They were part of the communication system.

My EDU486 Invisible Invaders camp work gives me another concrete picture of this principle. Youth could look, move, choose, handle tools, photograph, draw, write, explain, partner, pause, pass, or return. Identity Beads represented complementary contributions rather than ranking children by one performance. During evidence-to-action work and the showcase, youth could revise models, create advocacy messages, demonstrate, point, answer, ask, hold evidence, partner, or observe \citep{campEvidence2026}. Those routes were not all AAC, and I am not claiming that a particular camper used an AAC system. They do show how classroom design decides which forms of authorship become visible. A communication-rich classroom distributes ways to enter the work before a student is treated as disengaged, incapable, or unwilling.

\begin{figure}[!ht]
\centering
\newcommand{\agencybox}[2]{%
  \fcolorbox{PrimaryRed}{SoftBlue}{%
    \parbox[c][0.34in][c]{1.45in}{\centering\textbf{#1}\\[-0.01in]\small #2}%
  }%
}
\agencybox{Choose}{site, tool, and role}
\hspace{0.01in}{\color{PrimaryRed}$\longrightarrow$}\hspace{0.01in}
\agencybox{Contribute}{question, model, or message}
\hspace{0.01in}{\color{PrimaryRed}$\longrightarrow$}\hspace{0.01in}
\agencybox{Advocate}{vote, decline, or teach}
\hspace{0.01in}{\color{PrimaryRed}$\longrightarrow$}\hspace{0.01in}
\agencybox{Redesign}{time, access, and sequence}
\caption{Adapted from my Module 2, \textit{Who Has More Power After Support?} Invisible Invaders made agency observable through youth decisions and contributions, not by appearance, speech speed, or quiet compliance.}
\label{fig:camp-agency}
\end{figure}

My Module 3 paper, \textit{Communication That Travels With the Student: A Neurodiversity-Affirming, Culturally and Health Responsive Support Plan}, turned those lessons into a prospective plan for three fictional composite students \citep{garciabautista2026communicationplan}. Mateo made bilingual and culturally sustaining access visible; Eli made speech and multimodal access visible; and Luca made chronic illness, body state, privacy, low speech, and speech-related trauma visible. Their profiles are fictional, while the support strategies are grounded in de-identified camp and placement evidence. Together, the profiles kept me from designing for an abstract ``AAC student.'' A dependable system must survive changes in language, energy, pain, partner, tool, setting, routine, and willingness to speak.

\begin{tikzpicture}[
  piece/.style={
    draw=DeepNavy,
    rounded corners=2pt,
    text width=4.10in,
    align=center,
    font=\normalsize,
    inner xsep=-65pt,
    inner ysep=1pt
  },
  topwide/.style={
    draw=DeepNavy,
    rounded corners=2pt,
    text width=4.10in,
    align=center,
    font=\normalsize,
    inner xsep=-65pt,
    inner ysep=1pt
  },
  bottomwide/.style={
    draw=DeepNavy,
    rounded corners=2pt,
    text width=9.10in,
    align=center,
    font=\normalsize,
    inner xsep=-65pt,
    inner ysep=1pt
  },
  centerpiece/.style={
    draw=PrimaryRed,
    line width=0.9pt,
    rounded corners=2pt,
    fill=SoftGold,
    text width=2.15in,
    align=center,
    font=\normalsize\bfseries,
    inner xsep=7pt,
    inner ysep=1pt
  },
  flow/.style={
    -{Latex[length=2.1mm]},
    line width=0.75pt,
    draw=PrimaryRed
  }
]

% Center
\node[centerpiece] (center)
  {Reliable student authorship};

% Middle left/right
\node[piece, fill=SoftGray,
  left=0.16in of center]
  (partner) {Partner practice ask, wait, protect};

\node[piece, fill=SoftGold,
  right=0.16in of center]
  (family) {Family + home what travels};

% Full-width top left/right
\node[topwide, fill=SoftBlue,
  above=0.18in of partner]
  (mateo) {Mateo language + culture};

\node[topwide, fill=SoftGreen,
  above=0.18in of family]
  (eli) {Eli speech + multimodal access};

% Full combined-width bottom
\node[bottomwide, fill=SoftRose,
  below=0.28in of center]
  (luca) {Luca health + speech trauma};

% Connections
\draw[flow] (mateo.south east) -- (center.north west);
\draw[flow] (eli.south west) -- (center.north east);
\draw[flow] (partner.east) -- (center.west);
\draw[flow] (family.west) -- (center.east);
\draw[flow] (luca.north) -- (center.south);

\end{tikzpicture}

\Needspace{5\baselineskip}
That plan also forced adult behavior into the intervention. Premature suggestions can replace a student's idea; taking over work can erase authorship; rigid pacing can make fatigue look like a learner deficit; and visible adult dysregulation can make communication unsafe. Partner training must therefore include asking before helping, keeping hands off student work without permission, waiting, noticing all reliable modes, returning control after assistance, and repairing when adult behavior restricts access.
